\documentclass[12pt]{article}
\usepackage[utf8]{inputenc}
\usepackage[T1]{fontenc}
\usepackage[margin=1in]{geometry}
\usepackage{setspace}
\usepackage{xcolor}
\usepackage[colorlinks=true, urlcolor=blue!60!black, linkcolor=blue!60!black]{hyperref}
\setlength{\parindent}{0em}
\setlength{\parskip}{0.45em}
\renewcommand{\baselinestretch}{1.0}
\begin{document}
\begin{center}
{\Large \textit{City of Cleburne v. Cleburne Living Center, Inc.}, 473 U.S. 432 (1985)\par}
\vspace{0.3em}
{\normalsize Case Brief --- One-Page Edition\par}
\end{center}

\textbf{Memory Jogger.} The group-home case: rational basis review for intellectual disability, applied ``with bite'' to strike down prejudice-driven zoning (\href{https://supreme.justia.com/cases/federal/us/473/432/}{473 U.S. 432}).

\textbf{Facts.} The city required a special use permit for a regulated group home for thirteen adults with intellectual disabilities, treating it as a ``hospital for the feebleminded,'' though apartments, dormitories, and nursing homes were allowed by right. The Council denied the permit, citing neighbors' fears, flood risk, liability, and size.

\textbf{Procedural History.} The district court upheld the ordinance; the Fifth Circuit held the class quasi-suspect and reversed.

\textbf{Votes.} Affirmed in part, vacated in part: 9--0 on the judgment, 6--3 on the standard; White for the Court, Stevens concurring, Marshall concurring in part and dissenting in part.

\textbf{Holding.} Intellectual disability is not a quasi-suspect class, but ``mere negative attitudes, or fear'' and ``irrational prejudice'' are not legitimate interests, so the permit requirement is invalid as applied.

\textbf{Reasoning.} The group's differences are relevant to legitimate goals, and remedial legislation shows political responsiveness. But the city's reasons applied equally to freely permitted uses, leaving only prejudice. Stevens saw ``a continuum of judgmental responses''; Marshall called the method ``second order'' rational basis review and would have applied heightened scrutiny.

\textbf{Examples.} \emph{Same-Side:} a permit only for a group home for adults with autism is invalid. \emph{Opposite-Side:} cognitive testing for bus drivers is upheld. \emph{Fence-Sitter:} a permit rule for all large care facilities that exempts dormitories.

\textbf{Critique.} Scholars call it ``rational basis with bite,'' searching review under a deferential label; antisubordination critics fault denying heightened scrutiny because legislatures had helped the group.

\textbf{Key Quotation.} ``The Equal Protection Clause of the Fourteenth Amendment commands that no State shall `deny to any person within its jurisdiction the equal protection of the laws,' which is essentially a direction that all persons similarly situated should be treated alike.'' --- Justice White, for the Court, 473 U.S. at 439.

\end{document}
